% !TEX root = ../main.tex
\chapter{Fundamentals: Fixed Points, Linearization, Invariant Manifolds, Bifurcations, Chaos}\label{ch:fundamentals}
\begin{paracol}{2}

\begin{sources}
\begin{tabularx}{\linewidth}{@{}V l L@{}}
\toprule
\textbf{Video} & \textbf{Length} & \textbf{Content}\\
\midrule
\seg{2}{1}{V2.1 Topics in Dynamical Systems: Fixed Points, Linearization, Invariant Manifolds, Bifurcations \& Chaos} & 32:11 & Linear, mildly nonlinear and chaotic systems; fixed points and linearization; eigenvector coordinates; the Duffing oscillator; stable and unstable manifolds; pitchfork and Hopf bifurcations; discrete-time maps and the logistic map; the flow map; particles in flows and FTLE; chaos as stretching and folding.\\
\bottomrule
\end{tabularx}

\smallskip
\textbf{Transcript} \file{materials/transcripts/C02-V01.txt}.
\textbf{Book} Sec.~7.1 (``Dynamical systems'', ``Discrete-time systems'', ``Linear dynamics and
spectral decomposition''), Sec.~8.2 (linear systems).
This is the \emph{last} video of the playlist (recorded about three years after most of the others):
it is placed here because it is the background for everything that follows.
\end{sources}

\section*{Motivation}
A ``mile-high'' tour of what one learns in a course on dynamical systems \ts{2}{1}{0:20}. Every
vignette (fixed points, linearization, manifolds, bifurcations, maps, chaos) would be a couple of
lectures in a real course \ts{2}{1}{31:11}. The reason it belongs in a course on machine learning is
stated at the very end: when a model of a dynamical system is learned from data, these are the
properties one checks it against. Are the fixed points in the right places? Do the local linearizations
have the right stability? Does the model mix statistically like the data? \ts{2}{1}{31:32}.

For a physics MSc most of this is a refresher of classical mechanics and of the theory of ODEs; the
chapter is kept compact, with the derivations that are reused later (linearization, eigenvector
coordinates, flow map) written out in full.

% ------------------------------------------------------------------
\section{From linear to chaotic}\label{sec:fu-progression}

A dynamical system is any system that evolves in time according to some rules
\ts{2}{1}{0:41}: an apple falling under $F=ma$, a fluid obeying the Navier--Stokes equations, a
pendulum, a car, but also the brain, whose neurons fire according to rules that certainly exist even if
we cannot write them all down \ts{2}{1}{1:02}\ts{2}{1}{1:22}. Brunton orders systems by increasing
nonlinearity \ts{2}{1}{1:42}:
\begin{enumerate}
  \item \textbf{linear} systems: the pendulum at small angles, a network of springs, masses and dampers;
        they oscillate, and superposition holds \ts{2}{1}{2:03}\ts{2}{1}{2:23};
  \item \textbf{mildly nonlinear} systems, such as Hamiltonian or Lagrangian dynamics: the Duffing
        oscillator, a bead sliding on a wire shaped like a double well, with several fixed points
        \ts{2}{1}{3:05}\ts{2}{1}{3:26};
  \item \textbf{chaotic} systems: still deterministic, but with limited forecast ability, because an
        imperfectly known initial state or dynamics makes the prediction diverge from the truth after
        some time (double pendulum, Lorenz attractor, weather, turbulence) \ts{2}{1}{3:46}\ts{2}{1}{4:06}.
\end{enumerate}

\begin{definition}[New concepts: linearity, superposition]
The system $\dot{\vb x}=\vb f(\vb x)$ is \textbf{linear} if $\vb f(\vb x)=A\vb x$ for a matrix $A$. Then
\textbf{superposition} holds: if $\vb x_1(t)$ and $\vb x_2(t)$ are solutions, so is
$a\vb x_1(t)+b\vb x_2(t)$; adding two initial conditions adds the two solutions \ts{2}{1}{2:44}.
Superposition holds \emph{only} for linear systems.
\end{definition}

The state is again the minimal description, $(\theta,\dot\theta)$ for the pendulum, and $\vb f$ comes
from the Euler--Lagrange equations, $F=ma$ or Hamilton's equations \ts{2}{1}{4:48}\ts{2}{1}{5:09}
(\cref{def:ov-anatomy}).

% ------------------------------------------------------------------
\section{Fixed points and linearization}\label{sec:fu-linearization}

The pendulum behaves linearly when kicked a little near the bottom, nonlinearly when released from
near $\theta=\pi$ \ts{2}{1}{5:29}\ts{2}{1}{5:50}. The standard strategy for a nonlinear system is:
find the fixed points, linearize around them, and do linear analysis there \ts{2}{1}{9:35}.

\begin{definition}[New concepts: fixed point, Jacobian, linearization]\label{def:fu-fixed}
\begin{enumerate}[(a)]
  \item A \textbf{fixed point} (equilibrium) is a state $\bar{\vb x}$ with $\vb f(\bar{\vb x})=0$: started
        there, the system stays there \ts{2}{1}{6:10}.
  \item The \textbf{Jacobian} of $\vb f$ at $\bar{\vb x}$ is the matrix of partial derivatives
        $A=D\vb f(\bar{\vb x})$, $A_{ij}=\partial f_i/\partial x_j\,(\bar{\vb x})$ \ts{2}{1}{8:54}.
  \item The \textbf{linearization} at $\bar{\vb x}$ is the linear system
        $\frac{\mathrm d}{\mathrm dt}\Delta\vb x=A\,\Delta\vb x$ for the displacement
        $\Delta\vb x=\vb x-\bar{\vb x}$.
\end{enumerate}
The name \emph{Jacobian} honours Carl Gustav Jacob Jacobi (1804--1851), who studied these
determinants of partial derivatives in 1841.
\end{definition}

The pendulum has two fixed points: down, $(\theta,\dot\theta)=(0,0)$, and up, $(\pi,0)$; the second is
unstable, but perfectly balanced it stays put \ts{2}{1}{6:10}\ts{2}{1}{6:31}
(\cref{fig:ov-pendulum}).

\paragraph{Derivation.} Zoom into a neighbourhood of $\bar{\vb x}$ and write
$\vb x=\bar{\vb x}+\Delta\vb x$ with $\Delta\vb x$ small \ts{2}{1}{7:12}. Since $\bar{\vb x}$ is
constant,
\begin{equation}\label{eq:fu-taylor}
  \frac{\mathrm d}{\mathrm dt}\Delta\vb x=\dot{\vb x}=\vb f(\bar{\vb x}+\Delta\vb x)
  =\underbrace{\vb f(\bar{\vb x})}_{=0}+D\vb f(\bar{\vb x})\,\Delta\vb x+O(\lVert\Delta\vb x\rVert^2)
\end{equation}
by Taylor expansion \ts{2}{1}{7:32}\ts{2}{1}{7:52}. Close to the fixed point the higher-order terms are
negligible, and what remains is an honest linear system $\Delta\dot{\vb x}=A\Delta\vb x$
\ts{2}{1}{8:13}\ts{2}{1}{9:15}. Its eigenvalues say whether the fixed point is stable; its eigenvectors
say along which directions it is unstable \ts{2}{1}{9:15}.

\begin{theorem}[Hartman--Grobman]\label{thm:fu-hg}
If no eigenvalue of $A=D\vb f(\bar{\vb x})$ has zero real part (\emph{hyperbolic} fixed point), then in
a neighbourhood of $\bar{\vb x}$ the flow of $\dot{\vb x}=\vb f(\vb x)$ is topologically conjugate
(by a homeomorphism) to the flow of $\Delta\dot{\vb x}=A\Delta\vb x$.
\end{theorem}
\Cref{thm:fu-hg} is what justifies the strategy, and its hypothesis is what limits it: for eigenvalues
on the imaginary axis (centres, as in the Duffing wells below) the linearization alone decides nothing.

\paragraph{Eigenvector coordinates.} As in \cref{prop:ov-linear}, the solution is
$\vb x(t)=e^{At}\vb x_0$, which advances any initial condition any time forward or backward, and is
easy to compute \ts{2}{1}{9:55}\ts{2}{1}{10:16}. With the eigenvectors $T$ of $A$ and $\vb x=T\vb z$:
$\dot{\vb z}=T^{-1}\dot{\vb x}=T^{-1}A\vb x=T^{-1}AT\,\vb z=\Lambda\vb z$; it looks worse, but
$T^{-1}AT=\Lambda$ is diagonal, so the dynamics decouple \ts{2}{1}{10:37}\ts{2}{1}{10:57}\ts{2}{1}{11:18}.

\begin{intuition}[Three steps: map in, advance, map out]
Brunton's favourite reading of $\vb x(t)=Te^{\Lambda t}T^{-1}\vb x_0$ \ts{2}{1}{11:59}, read from
right to left \ts{2}{1}{12:40}:
\begin{enumerate}
  \item $T^{-1}$ maps the initial condition into eigenvector coordinates, $\vb z_0=T^{-1}\vb x_0$;
  \item there the system is decoupled: each $z_j$ is advanced independently, $z_j(t)=e^{\lambda_jt}z_j(0)$;
  \item $T$ maps back to physical space \ts{2}{1}{13:00}\ts{2}{1}{13:41}.
\end{enumerate}
The same ``lift, evolve linearly, project back'' structure returns in DMD and Koopman theory, where the
lift is nonlinear.
\end{intuition}

\begin{exercise}[Pendulum fixed points]
Linearize $\dot\theta=\omega$, $\dot\omega=-\sin\theta$ at $(0,0)$ and $(\pi,0)$, compute eigenvalues
and eigenvectors, and compare with \cref{fig:ov-pendulum}. Why does Hartman--Grobman say nothing about
the stability of $(0,0)$ by itself? What does energy conservation add?
\end{exercise}

% ------------------------------------------------------------------
\section{The Duffing oscillator and invariant manifolds}\label{sec:fu-duffing}

\begin{casestudy}[Particle in a double well]
The Duffing oscillator without friction \ts{2}{1}{14:43}:
\begin{equation}\label{eq:fu-duffing}
  \dot x=v,\qquad \dot v=x-x^3 .
\end{equation}
It is Newton's law $\ddot x=-V'(x)$ in the double-well potential $V(x)=x^4/4-x^2/2$ (a frictionless
roller coaster with this shape) \ts{2}{1}{17:08}\ts{2}{1}{17:28}. In the video the potential is written
loosely as ``$x^4-x^2$'', which has the same shape.
\begin{itemize}
  \item \emph{Fixed points}: $v=0$ and $x-x^3=0$, i.e.\ $x\in\{0,\pm1\}$ \ts{2}{1}{15:04}.
  \item \emph{Phase portrait}: a particle released at rest on a wall accelerates to maximal speed at the
        bottom, decelerates to rest on the opposite wall, comes back: a closed orbit in the $(x,v)$
        plane, forever, without friction \ts{2}{1}{15:44}\ts{2}{1}{16:27}\ts{2}{1}{16:47}.
  \item \emph{Jacobian} \ts{2}{1}{17:49}:
        $A(\bar x)=\begin{pmatrix}0&1\\1-3\bar x^2&0\end{pmatrix}$; friction would add a term in the
        lower right corner \ts{2}{1}{18:10}.
  \item At $(\pm1,0)$: $A=\begin{pmatrix}0&1\\-2&0\end{pmatrix}$, eigenvalues $\pm i\sqrt2$: neutrally
        stable, oscillation forever (a \emph{centre}) \ts{2}{1}{18:30}.
  \item At $(0,0)$: $A=\begin{pmatrix}0&1\\1&0\end{pmatrix}$, eigenvalues $\pm1$: a \emph{saddle};
        a small push sends the car away \ts{2}{1}{18:51}. The eigenvectors $(1,\pm1)/\sqrt2$ are the
        45-degree lines of \cref{fig:fu-duffing} \ts{2}{1}{19:11}\ts{2}{1}{19:31}.
\end{itemize}
\end{casestudy}

\begin{code}[fixed points and stability of the Duffing oscillator]
import numpy as np

def f(x):                      # Duffing without friction: x' = v, v' = x - x^3
    return np.array([x[1], x[0] - x[0]**3])

def jacobian(x):
    return np.array([[0.0, 1.0],
                     [1.0 - 3.0 * x[0]**2, 0.0]])

for xbar in ([0.0, 0.0], [1.0, 0.0], [-1.0, 0.0]):
    xbar = np.array(xbar)
    assert np.allclose(f(xbar), 0)           # it is a fixed point
    lam, T = np.linalg.eig(jacobian(xbar))
    print(xbar, 'eigenvalues', np.round(lam, 3))
    if np.isreal(lam).all():
        print('   eigenvectors (columns)\n', np.round(T, 3))
# [0. 0.]  eigenvalues [ 1.+0.j -1.+0.j]   eigenvectors (0.707, 0.707), (-0.707, 0.707)
# [1. 0.]  eigenvalues [0.+1.414j 0.-1.414j]
# [-1. 0.] eigenvalues [0.+1.414j 0.-1.414j]
\end{code}

\begin{nfigure}
\begin{tikzpicture}[x=2.2cm,y=2.2cm,>=Stealth]
  % level sets of E = v^2/2 - x^2/2 + x^4/4
  \foreach \c in {-0.2,-0.1} {
    \pgfmathsetmacro{\a}{sqrt(1-sqrt(1+4*\c))}\pgfmathsetmacro{\b}{sqrt(1+sqrt(1+4*\c))}
    \foreach \s in {1,-1} {
      \draw[intcol,thick,smooth,samples=60,domain=\a+0.0005:\b-0.0005] plot ({\s*\x},{sqrt(max(0,2*(\c+\x*\x/2-\x*\x*\x*\x/4)))});
      \draw[intcol,thick,smooth,samples=60,domain=\a+0.0005:\b-0.0005] plot ({\s*\x},{-sqrt(max(0,2*(\c+\x*\x/2-\x*\x*\x*\x/4)))});
    }
  }
  % separatrix E = 0 (homoclinic orbits)
  \draw[thmcol,very thick,smooth,samples=80,domain=-1.4142:1.4142] plot (\x,{abs(\x)*sqrt(max(0,1-\x*\x/2))});
  \draw[thmcol,very thick,smooth,samples=80,domain=-1.4142:1.4142] plot (\x,{-abs(\x)*sqrt(max(0,1-\x*\x/2))});
  % outer orbit E = 0.3
  \draw[defcol,thick,smooth,samples=160,domain=-1.5757:1.5757] plot (\x,{sqrt(max(0,2*(0.3+\x*\x/2-\x*\x*\x*\x/4)))});
  \draw[defcol,thick,smooth,samples=160,domain=-1.5757:1.5757] plot (\x,{-sqrt(max(0,2*(0.3+\x*\x/2-\x*\x*\x*\x/4)))});
  % eigenvectors at the saddle
  \draw[dashed,thmcol] (-0.5,-0.5) -- (0.5,0.5) (-0.5,0.5) -- (0.5,-0.5);
  \draw[->] (-1.8,0) -- (1.8,0) node[right] {$x$};
  \draw[->] (0,-1.05) -- (0,1.1) node[above] {$v$};
  \fill[thmcol] (0,0) circle (1.5pt); \fill[intcol] (1,0) circle (1.5pt) (-1,0) circle (1.5pt);
  \node[below,font=\scriptsize] at (1,0) {$1$}; \node[below,font=\scriptsize] at (-1,0) {$-1$};
  % potential, drawn above
  \begin{scope}[yshift=-2.6cm]
  \end{scope}
\end{tikzpicture}\hspace{0.8cm}
\begin{tikzpicture}[x=1.4cm,y=4cm]
  \draw[->] (-1.7,0) -- (1.7,0) node[right] {$x$};
  \draw[->] (0,-0.3) -- (0,0.45) node[above] {$V(x)$};
  \draw[srccol,very thick,smooth,samples=80,domain=-1.6:1.6] plot (\x,{\x*\x*\x*\x/4-\x*\x/2});
  \fill[intcol] (1,-0.25) circle (2pt) (-1,-0.25) circle (2pt);
  \fill[thmcol] (0,0) circle (2pt);
\end{tikzpicture}
\caption{Left: phase portrait of the Duffing oscillator \eqref{eq:fu-duffing}, level sets of the energy
$E=\tfrac12v^2+V(x)$. Green: oscillations inside one well (around the centres $(\pm1,0)$); red: the two
homoclinic orbits through the saddle, which are at the same time its stable and unstable manifolds;
blue: an orbit above the barrier, visiting both wells. Dashed: the eigenvectors of the saddle, tangent
to the manifolds. Right: the potential $V(x)=x^4/4-x^2/2$.}\label{fig:fu-duffing}
\end{nfigure}

\paragraph{Invariant manifolds.} Continuing the eigenvector lines of the saddle beyond the linear
neighbourhood sweeps out the \emph{stable} and \emph{unstable manifolds} \ts{2}{1}{19:52}. A second
example has a single fixed point at the origin which is a saddle, stable along $y$ and unstable along
$x$: an initial condition first approaches along the stable direction, then leaves along the unstable
one \ts{2}{1}{19:52}\ts{2}{1}{20:12}.

\begin{definition}[New concepts: stable, unstable, invariant manifolds]
For a fixed point $\bar{\vb x}$:
\begin{enumerate}[(a)]
  \item the \textbf{stable manifold} is $W^s=\{\vb x_0:\ \vb x(t)\to\bar{\vb x}\text{ as }t\to+\infty\}$;
  \item the \textbf{unstable manifold} is $W^u=\{\vb x_0:\ \vb x(t)\to\bar{\vb x}\text{ as }t\to-\infty\}$;
  \item a set is \textbf{invariant} if trajectories starting in it stay in it for all times; $W^s$ and
        $W^u$ are invariant.
\end{enumerate}
By the stable manifold theorem, at a hyperbolic fixed point $W^s$ and $W^u$ are smooth manifolds tangent
at $\bar{\vb x}$ to the stable and unstable eigenspaces of the linearization. \emph{Manifold}, the
English rendering of Riemann's \emph{Mannigfaltigkeit}, ``multiplicity'': a space that looks locally
like $\R^k$.
\end{definition}

\begin{example}[A saddle with a curved manifold]\label{ex:fu-saddle}
A standard system with one saddle (the book uses it for Koopman theory in Sec.~7.4) is
$\dot x=\mu x$, $\dot y=\lambda(y-x^2)$ with $\lambda<0<\mu$. The $y$-axis $x=0$ is invariant and is
the stable manifold. The unstable manifold is the parabola $y=\frac{\lambda}{\lambda-2\mu}x^2$:
substituting $y=cx^2$ gives $\dot y=2cx\dot x=2c\mu x^2$, which equals $\lambda(c-1)x^2$ for
$c=\lambda/(\lambda-2\mu)$. Its tangent at the origin is the $x$-axis, the unstable eigenvector.
\end{example}

\begin{intuition}[Saddles amplify uncertainty]
A ball of uncertainty advected near a saddle is squashed along the stable direction and stretched
along the unstable one \ts{2}{1}{20:33}\ts{2}{1}{20:54}. Repeat this with saddles whose manifolds
fold back (as in a bounded attractor) and you get \emph{stretching and folding}: one way to think
about chaos.
\end{intuition}

\begin{exercise}[Damped Duffing]
Add friction: $\dot v=x-x^3-\delta v$, $\delta>0$.
\begin{enumerate}[(a)]
  \item Compute the Jacobian at the three fixed points and classify them as a function of $\delta$
        (stable focus or node; saddle).
  \item Show that $E=\tfrac12v^2+V(x)$ satisfies $\dot E=-\delta v^2\le0$, and sketch the phase
        portrait. What are the basins of attraction of the two wells?
\end{enumerate}
\end{exercise}

% ------------------------------------------------------------------
\section{Bifurcations}\label{sec:fu-bifurcations}

Systems often depend on a parameter $\mu$: the inflow velocity or the size of the body in a wind
tunnel, the infection rate of a disease like COVID-19 \ts{2}{1}{21:15}\ts{2}{1}{21:35}. One wants to
know how the dynamics changes with $\mu$.

\begin{definition}[New concepts: bifurcation, bifurcation diagram]
A \textbf{bifurcation} is a qualitative change of the dynamics (number or stability of fixed points,
appearance of periodic orbits, \dots) at a parameter value $\mu_c$: a tiny change of $\mu$ produces a
big change of behaviour \ts{2}{1}{25:00}. A \textbf{bifurcation diagram} plots the fixed points (or
attractors) against $\mu$. From the Latin \emph{bifurcus}, ``two-pronged fork''.
\end{definition}

\paragraph{Pitchfork.} $\dot x=\mu x-x^3=x(\mu-x^2)$ \ts{2}{1}{21:56}. For $\mu<0$ everything
contracts to $x=0$. For $\mu>0$ the origin is linearly unstable: $x$ grows exponentially until the
cubic term saturates it at $x^2=\mu$ \ts{2}{1}{22:16}\ts{2}{1}{22:57}\ts{2}{1}{23:18}. The linearization
$\partial_x(\mu x-x^3)=\mu-3x^2$ gives $\mu$ at the origin and $-2\mu$ at $x=\pm\sqrt\mu$: two new stable
branches \ts{2}{1}{23:38}.

\paragraph{Hopf.} The same mechanism in two dimensions \ts{2}{1}{23:58}: in polar coordinates
\begin{equation}\label{eq:fu-hopf}
  \dot r=\mu r-r^3,\qquad \dot\theta=\omega ,
\end{equation}
i.e.\ $\dot x=\mu x-\omega y-x(x^2+y^2)$, $\dot y=\omega x+\mu y-y(x^2+y^2)$. For $\mu<0$ the origin
is a stable spiral; at $\mu=0$ a pair of complex eigenvalues $\mu\pm i\omega$ crosses the imaginary
axis, and for $\mu>0$ trajectories settle on a stable \emph{limit cycle} of radius $\sqrt\mu$, where the
linear growth balances the cubic term \ts{2}{1}{24:19}. In polar form it is ``directly analogous to the
pitchfork'' \ts{2}{1}{24:40}. Hopf bifurcations are everywhere in physical instabilities, for example
the onset of vortex shedding behind a cylinder.

\begin{nfigure}
\begin{tikzpicture}[x=1.6cm,y=1.4cm,>=Stealth]
  \draw[->] (-1.6,0) -- (1.7,0) node[right] {$\mu$};
  \draw[->] (0,-1.4) -- (0,1.5) node[above] {$\bar x$};
  \draw[intcol,very thick] (-1.5,0) -- (0,0);
  \draw[thmcol,very thick,dashed] (0,0) -- (1.5,0);
  \draw[intcol,very thick,smooth,domain=0:1.5,samples=60] plot (\x,{sqrt(\x)});
  \draw[intcol,very thick,smooth,domain=0:1.5,samples=60] plot (\x,{-sqrt(\x)});
  \foreach \y in {0.9,-0.9} \draw[->,srccol] (-0.8,\y) -- (-0.8,{\y*0.35});
  \draw[->,srccol] (0.8,0.15) -- (0.8,0.65); \draw[->,srccol] (0.8,-0.15) -- (0.8,-0.65);
  \draw[->,srccol] (0.8,1.35) -- (0.8,1.0);  \draw[->,srccol] (0.8,-1.35) -- (0.8,-1.0);
  \node[font=\footnotesize] at (1.25,1.45) {$\bar x=\sqrt\mu$};
\end{tikzpicture}\hspace{1.2cm}
\begin{tikzpicture}[scale=1.0,>=Stealth]
  \draw[->] (-1.7,0) -- (1.7,0) node[right] {$x$};
  \draw[->] (0,-1.7) -- (0,1.7) node[above] {$y$};
  \draw[intcol,very thick] (0,0) circle (1.1);
  \draw[defcol,thick,->,domain=0:1080,samples=200,smooth] plot ({(0.12+0.98*(1-exp(-\x/300)))*cos(\x)},{(0.12+0.98*(1-exp(-\x/300)))*sin(\x)});
  \draw[thmcol,thick,->,domain=0:900,samples=200,smooth] plot ({(1.6-0.48*(1-exp(-\x/250)))*cos(\x+40)},{(1.6-0.48*(1-exp(-\x/250)))*sin(\x+40)});
  \node[font=\footnotesize,intcol] at (1.2,-1.35) {$r=\sqrt\mu$};
\end{tikzpicture}
\caption{Left: pitchfork bifurcation of $\dot x=\mu x-x^3$ (solid: stable, dashed: unstable fixed
points; arrows: the flow in $x$ at fixed $\mu$). Right: Hopf normal form \eqref{eq:fu-hopf} for
$\mu>0$: trajectories from inside and outside spiral onto the limit cycle.}\label{fig:fu-bif}
\end{nfigure}

\begin{remark}
This is a cartoon of what is a whole graduate course: the classification of bifurcations (saddle-node,
transcritical, pitchfork, Hopf, \dots) and normal forms \ts{2}{1}{24:40}.
\end{remark}

\begin{exercise}[Hopf in Cartesian coordinates]
Starting from \eqref{eq:fu-hopf}, derive the equations for $x=r\cos\theta$, $y=r\sin\theta$, compute
the Jacobian at the origin, and check that its eigenvalues are $\mu\pm i\omega$. Solve the radial
equation exactly (hint: $u=1/r^2$) and show that every $r(0)>0$ tends to $\sqrt\mu$ for $\mu>0$.
\end{exercise}

\begin{exercise}[Subcritical pitchfork]
Study $\dot x=\mu x+x^3-x^5$. Draw the bifurcation diagram, find the saddle-node bifurcations, and
show the hysteresis loop obtained by slowly increasing and then decreasing $\mu$.
\end{exercise}

% ------------------------------------------------------------------
\section{Discrete-time systems and the flow map}\label{sec:fu-discrete}

\paragraph{Maps.} Dynamics can also be given in discrete time, $\vb x_{k+1}=\vb F(\vb x_k)$
\ts{2}{1}{25:21}. The \emph{logistic map} $x_{k+1}=\mu x_k(1-x_k)$ models a population, say of bunnies
in year $k$: compute the right-hand side to get next year's population, iterate year after year
\ts{2}{1}{25:41}\ts{2}{1}{26:02}. Here $\mu$ is a reproduction rate, and sweeping it gives very different
behaviours: around $\mu=3.45$ the population cycles reproducibly among a few values; for larger $\mu$ it
wanders over a fractal set, the discrete-time analogue of chaos \ts{2}{1}{26:22}\ts{2}{1}{26:43}.
Chapter~3 simulates it.

\paragraph{The flow map.} Continuous-time systems induce discrete-time ones, which is what a computer
simulation does \ts{2}{1}{27:04}\ts{2}{1}{27:24}.
\begin{definition}[New concepts: flow map]\label{def:fu-flowmap}
The \textbf{flow map} of $\dot{\vb x}=\vb f(\vb x)$ over a time $t$ is
\begin{equation}\label{eq:fu-flowmap}
  \vb F_t(\vb x(t_0))=\vb x(t_0+t)=\vb x(t_0)+\int_{t_0}^{t_0+t}\vb f(\vb x(\tau))\,\mathrm d\tau .
\end{equation}
The integration variable $\tau$ is there because the particle sees a different $\vb f$ as it moves
\ts{2}{1}{27:45}. Sampling with step $\Delta t$ gives the map $\vb x_{k+1}=\vb F_{\Delta t}(\vb x_k)$,
a program that steps the system forward \ts{2}{1}{28:06}. The flow maps form a group:
$\vb F_t\circ\vb F_s=\vb F_{t+s}$, $\vb F_0=\mathrm{id}$.
\end{definition}
For a linear system, $\vb F_t=e^{At}$ (\cref{prop:ov-linear}). The book adds that maps are more general
than flows: they include discontinuous and hybrid systems, and they are the natural setting for
sampled experimental data and digital control.

\paragraph{Particles in flows.} In fluid mechanics $\vb f$ is the velocity field and $\vb x$ the
position of a tracer particle; integrating many particles shows how the fluid mixes and how a buoy or a
pollutant spreads in the ocean \ts{2}{1}{28:26}\ts{2}{1}{29:08}. The field then depends on time, and the
\emph{finite-time Lyapunov exponent} (FTLE) fields are the time-varying analogues of stable and
unstable manifolds \ts{2}{1}{28:46}.

\begin{secondary}[Finite-time Lyapunov exponents]
For the flow map $\vb F_{t_0}^{t_0+T}$ of a time-dependent field, the Cauchy--Green tensor is
$C=(D\vb F)^{\mathsf T}D\vb F$ and the FTLE is
$\sigma_T(\vb x_0)=\frac{1}{|T|}\ln\sqrt{\lambda_{\max}(C)}$, the exponential rate at which nearby
particles separate. Ridges of $\sigma_T$ computed forward in time ($T>0$) mark repelling structures
(analogues of stable manifolds), backward ($T<0$) attracting ones (unstable manifolds): the Lagrangian
coherent structures of G.~Haller.
\end{secondary}

\begin{exercise}[Flow map of a linear system]
For $\dot{\vb x}=A\vb x$ with $A=\begin{pmatrix}-1&2\\-2&-1\end{pmatrix}$, compute $\vb F_t=e^{At}$ in
closed form. Sample at $\Delta t$: what are the eigenvalues of $\vb F_{\Delta t}$? Check numerically
(\texttt{scipy.linalg.expm}) and verify $\vb F_t\vb F_s=\vb F_{t+s}$.
\end{exercise}

% ------------------------------------------------------------------
\section{Chaos}\label{sec:fu-chaos}

Integrate a tight cube of initial conditions through the Lorenz system: they stay together for a
while, then stretch, fold and mix along the whole attractor \ts{2}{1}{29:28}\ts{2}{1}{29:49}\ts{2}{1}{30:09}.

\begin{definition}[New concepts: chaos, attractor, prediction horizon]
A deterministic system is \textbf{chaotic} if it has sensitive dependence on initial conditions: nearby
trajectories separate exponentially, at a rate given by the largest Lyapunov exponent $\lambda_1>0$, and
mix over a bounded \textbf{attractor} (the set the trajectories approach). An initial uncertainty
$\delta_0$ grows like $\delta_0e^{\lambda_1t}$, so a tolerance $\Delta$ is reached after the
\textbf{prediction horizon} $t_*\approx\lambda_1^{-1}\ln(\Delta/\delta_0)$; after that only statistical
statements are possible \ts{2}{1}{30:30}\ts{2}{1}{30:51}. From the Greek \emph{cháos}, the gaping void
before creation.
\end{definition}

The logarithm is the bad news: improving the initial measurement by a factor 10 only adds
$\lambda_1^{-1}\ln10$ to the horizon. Local weather (where a hurricane will go) and turbulence are
chaotic \ts{2}{1}{30:51}\ts{2}{1}{31:11}.

\begin{history}[Lorenz and the butterfly]
Edward Lorenz (MIT) found sensitive dependence in 1961 when he restarted a weather simulation from a
printout rounded to three decimals and the forecast diverged completely. His 1963 paper
\emph{Deterministic nonperiodic flow} (J.\ Atmos.\ Sci.\ 20, 130--141) reduced convection to the three
equations now named after him. The ``butterfly effect'' comes from the title of a 1972 talk at the AAAS:
\emph{Predictability: Does the flap of a butterfly's wings in Brazil set off a tornado in Texas?}\par
\textit{References:} E.~N.~Lorenz, \emph{The Essence of Chaos} (Univ.\ of Washington Press, 1993);
J.~Gleick, \emph{Chaos} (1987).
\end{history}

\begin{exercise}[Prediction horizon]
The largest Lyapunov exponent of the Lorenz system at standard parameters is $\lambda_1\approx0.906$.
With an initial error of $10^{-6}$ and a tolerance of $1$, estimate the horizon. How much does it grow
if the error is reduced to $10^{-12}$?
\end{exercise}

\begin{keyformulas}
\begin{itemize}
  \item Fixed point $\vb f(\bar{\vb x})=0$; linearization $\Delta\dot{\vb x}=D\vb f(\bar{\vb x})\Delta\vb x$;
        valid at hyperbolic points (Hartman--Grobman).
  \item $\vb x(t)=Te^{\Lambda t}T^{-1}\vb x_0$: map in, advance decoupled modes, map out.
  \item Duffing $\dot x=v$, $\dot v=x-x^3$: centres $(\pm1,0)$ with $\pm i\sqrt2$, saddle $(0,0)$ with $\pm1$.
  \item Pitchfork $\dot x=\mu x-x^3$, Hopf $\dot r=\mu r-r^3$, $\dot\theta=\omega$: branches/cycle at $\sqrt\mu$.
  \item Flow map $\vb F_t(\vb x_0)=\vb x_0+\int_0^t\vb f(\vb x(\tau))\dd\tau$; prediction horizon
        $t_*\approx\lambda_1^{-1}\ln(\Delta/\delta_0)$.
\end{itemize}
\end{keyformulas}

\end{paracol}
